% Generated draft: requires mathematical and source review.
\documentclass[11pt]{article}
\usepackage[margin=25mm]{geometry}
\usepackage{fontspec}
\setmainfont{Latin Modern Roman}
\usepackage{amsmath,amssymb,mathtools,unicode-math}
\setmathfont{Latin Modern Math}
\usepackage{enumitem,needspace,xurl,hyperref}
\setlength{\parindent}{0pt}
\setlength{\parskip}{5pt}
\setlength{\emergencystretch}{3em}
\newcommand{\N}{\mathbb{N}}
\newcommand{\Z}{\mathbb{Z}}
\newcommand{\Q}{\mathbb{Q}}
\newcommand{\R}{\mathbb{R}}
\newcommand{\C}{\mathbb{C}}
\newcommand{\F}{\mathbb{F}}
\newcommand{\A}{\mathbb{A}}
\newcommand{\PP}{\mathbb{P}}
\newcommand{\E}{\mathbb{E}}
\newcommand{\eps}{\varepsilon}
\newcommand{\dd}{\,\mathrm{d}}
\begin{document}
\section*{UnsolvedMath ID: 30005925}
\textbf{Source problem code:} catalogue entry
\subsection{Definitions and mathematical statement}
Write \(E\colon\mathcal A\to B\) for the \(B\)-valued expectation implicit in the source record, where \(z\in\mathcal A\). In the usual operator-valued probability terminology, such an expectation fixes \(B\), is unital, and satisfies \(E(b_1ab_2)=b_1E(a)b_2\) for \(a\in\mathcal A\) and \(b_1,b_2\in B\). A \(B\)-valued circular element is conventionally characterized by the vanishing of all operator-valued cumulants in \(z,z^*\), with intervening coefficients from \(B\), except possibly second-order cumulants containing one \(z\) and one \(z^*\). The supplied excerpt does not specify the ambient operator-algebraic assumptions or the precise source conventions for the decomposition terms below.

Define the moment span
\[
M_z:=\operatorname{span}\bigl\{E((z^*z)^k):k\in\mathbb Z_{\geq0}\bigr\}\subseteq B,
\]
where the zeroth power is the identity. Let \(\tau\) denote the topology intended by the source for the density condition; that topology is not identified in the supplied record.

\textbf{Question.} In the setting intended by the cited report, is it true that every \(B\)-valued circular element \(z\) which admits a free bipolar decomposition and satisfies
\[
\overline{M_z}^{\,\tau}=B
\]
admits a free polar decomposition?

\textbf{Unresolved qualifications.} The record does not define either a free bipolar decomposition or a free polar decomposition. In particular, it does not specify the factors permitted in these decompositions, the expectation or state relative to which freeness is required, whether freeness is amalgamated over \(B\), or whether a decomposition may be realized in an enlarged probability space. These conventions, the ambient assumptions on \(\mathcal A,B,E\), and the topology \(\tau\) must be recovered from the cited source before this question becomes a fully specified mathematical statement. No choice among them is asserted here.
\subsection{Short English statement}
The question asks whether a \(B\)-valued circular element with a free bipolar decomposition must also have a free polar decomposition when the linear span of its expected positive moments is dense in \(B\). The supplied excerpt leaves the decomposition conventions, ambient setting, and density topology unspecified, so this remains a qualified statement-only draft.
\subsection{Sources}
\begin{itemize}
\item Local UnsolvedMath registry, supplied complete source record: permanent ID \(30005925\), title \emph{Free Polar Decomposition Under a Nondegeneracy Condition}, source problem code \texttt{catalogue entry}, in the collection \emph{Oberwolfach Reports Open Problems}. The proposed year is recorded as 2024; no individual proposer is supplied.
\item \emph{Non-commutative Function Theory and Free Probability} (2024), Oberwolfach Reports, pp.~1227--1284, DOI 10.4171/owr/2024/22. Supplied report reference and literature-source link: \url{https://doi.org/10.4171/owr/2024/22}. The background describes this report as documenting workshop discussions and open problems.
\item Dykema--Griffin, supplied literature-source link: \url{https://arxiv.org/abs/2306.17133}. The record supplies no paper title or precise supporting location.
\item Literature assessment dated exactly \texttt{2026-08-22}: the supplied curation labels the problem open and states that degenerate examples and partial structural theorems are known, while existence under the stated dense-span nondegeneracy condition remains explicitly open. This is supplied dated curation, not an independently verified literature assessment; the record requests independent verification of time-sensitive status and cited claims.
\item The registry labels the statement \texttt{corrected\_verified} and reports that it was reconstructed and checked against the official Oberwolfach report. This is a supplied verification claim, not verification performed for this draft. The retained original excerpt explicitly states the dense-span condition but does not reproduce the full decomposition question.
\end{itemize}
\end{document}
